\begin{abstract}

本文围绕算力约束下大语言模型的资源配置问题，建立了一套从数据刻画、规律纳入、配置求解到历史检验的模型体系。四个子问题共用同一套数据口径与参数体系，前一问的输出构成后一问的输入。

针对问题一，本文对 22 个质量指标统一方向后合成质量评分。方向表按可信度分三层推得：题面或原文明示者可核验，字段定义与指标间一致性支持者可推证，两者皆不足者列为待定。按 token 求和的对数似然比属外延量，与文档长度的秩相关约为 $-0.93$，逐 token 归一化后其与教育价值打分的秩相关由 $-0.499$ 转为 $+0.374$，方向依据由建模判断升为数据可证。合成采用门控形式 $Q=V(1-P)$，使污染类指标以乘性惩罚进入而不与价值类指标在平均中相互抵消；冲突相应定义为同一文本在高价值与高污染两侧同时取高分，阈值按分位数锚定。在全部 272\,486 行质量信号上，门控合成使可分辨域对由 $19/21$ 增至 $21/21$，权重扰动下的域序 Kendall $\tau$ 由 $0.9595$ 升至 $0.9976$，二者差值区间不含零；域序最高一档为 $\{\text{arxiv},\text{c4},\text{commoncrawl}\}$，档内位次依语义分组而变，故不作断言。主冲突判据在 $p75$ 锚点下给出冲突率 $0.0408$（$p70$ 与 $p80$ 分别为 $0.0687$ 与 $0.0208$），其样本画像与其余判据近乎不相交。配比建模以去尺度偏移后的决定系数与秩相关为验证口径，主方法在 1m 尺度上优于线性基线 $+0.0703$，在 60m 与 1B 上分别低 $0.0050$ 与 $0.0464$，不构成跨尺度优势。引入域级 $Q$ 的去偏决定系数增量在三个尺度上均不显著，区间均含零；成因是 $Q$ 的构造本身含配比，二者存在构造性相关。这一零结果是对题目「是否以及如何引入 $Q$」的有证据回答。

针对问题二，本文把质量与配比并入数据量项，建立\linebreak $L=E_s+\allowbreak AN^{-\alpha}+\allowbreak B(DQ^{\gamma}\exp(\boldsymbol{b}^{\top}\boldsymbol{z}_{\boldsymbol{p}}))^{-\beta}$，在 $Q=1$ 且配比均匀时逐项退回经典标度律。数据条件决定估计顺序：主拟合轨迹无质量变化，含质量的半合成网格只有 5 个数据量取值，两者不能在同一份数据上同时辨识，故先由真实轨迹定 $\beta$、再由网格反解 $\gamma$。三项预声明判据全部通过（退化相对差 $1.01\times10^{-8}$；超越性降幅 $0.0419$，区间不含零；弹性符号两侧成立）。核心结论是质量指数 $\gamma=0.9659$，自助区间 $[0.8839,1.0555]$ 含 $1.0$，即质量的相对变化按比例放大有效 token 数；但该结论以「数据量指数可跨源共享」为前提，而该前提未通过检验（$\gamma$ 位移达其区间宽的 $1.246$ 倍）。若该指数改在含质量网格上自由估计，$\gamma=1.1798$，区间不含 $1.0$。两条路线本文一并报告。质量增加 $0.1$ 所补偿的参数量，在参照点 $N=1.0$B、$D=150.0$B、$Q=0.5$ 上为 $-0.1341$B，B7 侧为 $-0.0945$B，符号为负表示提质可减少参数量。

针对问题三，本文在预算约束 $kND+D[g(Q)-g(Q_0)]_+\le C$（$k=6+\eta L_{\text{ctx}}$）下最小化损失。内点条件不足以刻画最优解；本文按活跃约束集分支枚举，KKT 五项判据在三个预算锚点与 57 点扫描上全部通过。本文把结构性转移定义为三种并逐一识别，发现定义决定结论：活跃约束集变化给出位置 $\log_{10}C=21.4052$ 且不依赖阈值与网格步长；成本分项占比排序变化给出 $22.625$ 且依赖步长；分配份额突变在阈值 $0.05$ 下判不出，步长减半即由存在翻转为不存在。合读的结论是转移存在、形态为连续饱和而非突跳。成本函数的影响集中在预算紧张时：$C=10^{19}$ 档三类形式给出 $0.6460$、$0.5390$ 与 $1.0$ 三种最优质量，$10^{22}$ 与 $10^{24}$ 档完全一致。上下文长度跨越 $6/\eta=30000$ 时显著压缩可负担规模（$N^{*}$ 降 $54.8\%$、$D^{*}$ 降 $53.4\%$）而几乎不改变相对配比（$N{:}D$ 极差 $3.22\%$）。

针对问题四，本文把窗口内的能力提升拆为规模通道与规模之外的技术通道，并如实报告口径敏感性。稳健口径为趋势法：两通道月斜率之比为 $9.98$，规模份额 $0.0911$ 的自助区间含零。据此本文主张：窗口 2024 年 6 月至 2025 年 3 月内能力提升约 $5.4865$--$6.6592$ 分（四种端点/规格算法的极差），由非规模技术/数据工程综合进步主导，规模通道贡献与零不可区分。本文同时发现任何具体百分比分割都不稳定——端点定义变动使规模份额极差达 $0.2875$ 并发生符号翻转，故不给出百分比。同期高分段模型规模在收缩（参数量 0.9 分位由 $46.703$B 降至 $14.770$B），而前沿由 $25.9201$ 分升至 $41.3875$ 分。预测环节中性情景 12 个月中位数为 $56.330$ 分，但同口径的纯趋势基线给出 $60.290$ 分，主方法在点值与区间宽度上均不优于基线，本文同表列出。算力增长放缓对短期预测影响很小：情景极差随规模份额严格缩放，为 $11.786$、$4.1408$ 与 $0$ 分。

本文的创新在于：以门控形式解决方向相反指标共存时的合成问题，并由此自然导出冲突的定义；把不可检验的跨源共享前提显式化为条件，同时给出该前提成立与不成立两条读数；以三种定义的对照识别结构性转移，指出「存在」这一判断的成立范围；把点值的不稳定性归因到重抽样、函数规格与样本端点三个来源并分别扰动，据此撤回不可主张的百分比分割。

\keywords{标度律；数据质量；领域配比；算力约束优化；KKT 条件；结构性转移；前沿预测}
\end{abstract}
